\section{Năm nhóm vấn đề}
\label{sec:ch10-nam-nhom-van-de}

\index{LIME!năm nhóm vấn đề}%
Bảng thứ nhất trong ba bảng mà mục trước vừa kể, bảng phân loại chính, dựng
trên hai trục, và trục thứ nhất là câu hỏi một biến thể sinh ra để sửa cái gì. Khảo sát rút từ tài liệu ra năm nhóm vấn đề, và
mỗi nhóm được phát biểu kèm danh sách những bài báo đã nêu
nó~\autocite{p22whichlime}. Năm nhóm ấy như sau, theo thứ tự khảo sát đặt.

\index{tính cục bộ}%
Trước hết là \tn{tính cục bộ}{locality}. Lời giải thích có thể không đủ cục bộ
với điểm đang xét, tức không thật sự riêng cho điểm ấy, nếu các mẫu đã nhiễu dùng để dựng mô hình thay thế không đại
diện được cho biên quyết định ở ngay quanh điểm ấy. Bốn bài báo được dẫn cho
nhóm này~\autocite{p22whichlime}. Đây chính là câu hỏi mà mục~\ref{sec:ch03-dieu-lime-chua-giai-quyet}
để lại: một đường khớp tốt trên lân cận do người dùng dựng không tự chứng minh
lân cận ấy là lân cận đúng.

Nhóm thứ hai là độ khớp của mô hình thay thế. Khảo sát phát biểu nó là chuyện mô
hình thay thế \enquote{có thể không nắm bắt chính xác hành vi của mô hình gốc,
dẫn tới những lời giải thích không phản ánh đầy đủ quá trình ra quyết định của
mô hình gốc}, với sáu bài báo được dẫn~\autocite{p22whichlime}. Các tác giả
buộc nhóm này với nhóm trước:
độ khớp kém có thể sinh ra từ tính cục bộ kém, nhưng cũng có thể sinh ra từ một
mô hình thay thế quá đơn giản hoặc từ một biểu diễn
\tn{đặc trưng}{feature} không đủ.

Tên gọi của nhóm ấy cần một câu riêng, vì cùng một chữ tiếng Anh đang chỉ hai
thứ khác nhau trong hai bài báo mà sách đều dẫn. Khảo sát gọi nhóm vấn đề trên là
\enquote{fidelity}, và cái nó hỏi là đầu ra của
\tn{mô hình thay thế cục bộ}{local surrogate} có bám được đầu ra của mô hình gốc
trên vùng quanh điểm cần giải thích hay không. Sách gọi quan hệ ấy là độ khớp,
và giữ nó tách khỏi định nghĩa~\ref{def:ch01-do-trung-thuc}: độ khớp so hai đầu
ra trên một miền đã nói rõ, còn độ trung thực so lời giải thích với những phụ
thuộc của hành vi dưới một tập can thiệp đã nói rõ. Một mô hình thay thế khớp
tốt trên miền đã lấy mẫu vẫn có thể chỉ sai đặc trưng mà hành vi phụ thuộc vào,
nên độ khớp là bằng chứng đặt cạnh độ trung thực chứ không phải chính nó; đó là
hai dòng đầu của bảng~\ref{tab:ch01-bon-quan-he}. Đó
cũng không phải
đại lượng fidelity của
đánh giá \tn{phản thực}{counterfactual} mà chương~\ref{ch:khung-hop-nhat} đặt tên và
chương~\ref{ch:do-do-trung-thuc} đọc kỹ. Chữ fidelity vì vậy được giữ cho đại
lượng đã có tên ấy; chỗ duy nhất chương này viết chữ ấy nữa là khi nói về chính
chữ ấy.

Nhóm thứ ba là \tn{khả năng diễn giải}{interpretability}: cách trình bày lời
giải thích có thể không
diễn đạt được quyết định của mô hình theo lối người dùng đọc được, hoặc phải sửa
lại cho hợp một kiểu dữ liệu khác. Năm bài báo được
dẫn~\autocite{p22whichlime}. Nhóm này khác bốn
nhóm còn lại ở một điểm mà chương~\ref{ch:giai-thich-de-lam-gi} đã tách sẵn: nó nói về
người đọc lời giải thích, không nói về quan hệ giữa lời giải thích và mô hình.

\index{tính ổn định}%
Nhóm thứ tư là \tn{tính ổn định}{stability}, tức nhóm chứa trường hợp mà
chương~\ref{ch:chi-so-khong-do-do-trung-thuc} đã giao lại cho chương này. Khảo
sát phát biểu nó ở dạng phủ định, như một hỏng hóc: lời giải thích của LIME
\enquote{có thể thay đổi đáng kể do những thay đổi nhỏ ở dữ liệu đầu vào, ở quá
trình nhiễu, ở phép lấy mẫu, ở các lần chạy lặp lại, hoặc ở mô hình bên dưới},
và hệ quả là những kết quả không nhất quán và không đáng
tin~\autocite{p22whichlime}. Sách phát biểu lại điều kiện ấy ở dạng khẳng định
để về sau trích theo số; nội dung là của khảo sát, còn dạng khẳng định là của
sách.

\begin{definition}[Tính ổn định của lời giải thích]
\label{def:ch10-tinh-on-dinh}
Một phương pháp giải thích là ổn định khi những thay đổi nhỏ ở dữ liệu đầu vào,
ở \tn{phép nhiễu}{perturbation}, ở phép lấy mẫu, ở lần chạy lặp lại, hay ở mô
hình được giải thích chỉ làm lời giải thích thay đổi ít.
\end{definition}

Bảy bài báo được dẫn cho nhóm này, nhiều hơn mọi nhóm khác. Khảo sát nói thêm
một điều về hệ quả mà định nghĩa~\ref{def:ch10-tinh-on-dinh} chưa chứa: khi lời
giải thích đổi giữa hai lần chạy, người đọc \enquote{không biết được sự bất ổn
định bắt nguồn từ phần nào}, nên mất lòng tin vào phương pháp giải thích hoặc
vào chính mô hình được giải thích~\autocite{p22whichlime}. Đó đúng là chỗ
chương~\ref{ch:chi-so-khong-do-do-trung-thuc} dừng lại: hai lần chạy cho hai lời
giải thích khác nhau thì lời giải thích đang phụ thuộc vào một thứ nằm ngoài mô
hình và đầu vào, và kết luận ấy không cần ground truth nào. Cái nó không cho
biết là lần nào sai, hay có lần nào sai.

Nhóm thứ năm là hiệu năng: sinh các mẫu đã nhiễu, gọi mô hình trên từng mẫu rồi
khớp mô hình thay thế đều tốn thời gian. Ba bài báo được dẫn, ít nhất trong năm
nhóm~\autocite{p22whichlime}.

Năm nhóm ấy không độc lập, và khảo sát nói rõ điều đó bằng hai ví dụ ngược
chiều nhau: tăng tính cục bộ có thể làm hiệu năng xấu đi, còn chịu hiệu năng
kém đi, tức tốn thời gian hơn, lại có thể làm tính ổn định tốt
lên~\autocite{p22whichlime}. Nói cách khác, năm
nhóm là năm mặt của cùng một tập lựa chọn tự do, chứ không phải năm hỏng hóc
tách rời nhau, và mục~\ref{sec:ch10-vuon-bien-the} sẽ cho thấy hơn hai phần
năm số biến thể nhắm vào nhiều hơn một nhóm cùng lúc.
